\documentclass[11pt]{ctexart}
\usepackage[a4paper,margin=1in]{geometry}
\usepackage{graphicx}
\usepackage{xcolor}
\usepackage{hyperref}
\definecolor{refgray}{gray}{0.45}
\title{\textbf{市场最新观点汇总}\\[0.4em]{\large 地缘风险缓释与结构性分化并存：欧洲韧性、亚洲汇市重估、AI硬件瓶颈与港股定价权迁移}}
\author{KC桌面 自动生成}
\date{260525}
\begin{document}
\maketitle
\section*{一页摘要}
\begin{itemize}
  \item 市场对西班牙的结构性韧性存在系统性低估，其主权利差中几乎未计入国内风险溢价，利差走阔完全由全球因素驱动。
  \item 韩国央行即将进行有分寸的鹰派转向，首次正式承认结构性通胀，但科技与非科技部门的剪刀差扩大使政策决策复杂化。
  \item 美伊协议可能性上升，市场低估了地缘风险缓释后亚洲资产的系统性重估，人民币升值潜力被结构性低估。
  \item 人民币中间价模型显示逆周期因子正在收敛，定价机制回归基准逻辑，汇率双向波动可能成为可交易的真实状态。
  \item 恒生指数调整揭示被动资金正系统性地将港股定价重心从传统行业转向互联网平台、医疗健康和AI软件服务。
  \item 中国燃气行业面临成本端的不对称风险，工业气量回暖具有假性复苏特征，可持续性取决于外部扰动。
  \item 山西煤矿事故可能触发中央层面更严格的安全检查，市场低估了安全冲击对供给侧长期约束力的再定价能力。
  \item AI半导体的下一个瓶颈不在芯片本身而在先进封装，封装产能的稀缺性正在重新定义供给曲线和竞争格局。
  \item 香奈儿FY25财报显示品牌势能复苏，美国市场是核心引擎，亚太企稳是信心指标，资本开支转向垂直整合。
\end{itemize}
\section{宏观与利率：欧洲韧性重估与亚洲政策分化}
\textbf{欧洲核心国家（西班牙）的结构性转型被市场低估，而韩国央行面临科技与非科技分化的政策两难，全球宏观叙事正从单一通胀交易转向结构性分化交易。}

\begin{itemize}
  \item 西班牙的结构性韧性被系统性低估：其能源强度在EMU4中最低，财政选择性地将资源投入能源补贴而非国防开支，使得财政平衡处于2007年以来最佳水平，且是唯一预计未来三年债务GDP比下降的EMU4经济体。
  \item 西班牙主权利差中几乎未计入任何国内风险溢价，当前利差走阔完全由全球通胀和增长放缓驱动，若市场重新定价国内风险，利差存在显著压缩空间。
  \item 韩国央行即将在5月会议上进行有分寸的鹰派转向，从中性调整为审慎关注上行风险，标志着其首次正式将能源冲击和房价再升温同时纳入政策框架。
  \item 韩国经济呈现双轨复苏：科技生产环比增长11\%，非科技产出基本持平，复苏高度集中在占GDP约10\%的半导体行业，给央行带来加息抑或维持宽松的两难。
  \item 韩国通胀上修压力主要在能源端，核心通胀可控，工资增速持续放缓，半导体就业占比不到1\%，不会传导至整体劳动力市场。
\end{itemize}
\begin{figure}[htbp]
\centering
\includegraphics[width=0.88\linewidth]{figures/fig_001.jpg}
\caption{Exhibit 1 - Exhibit 1: Growth Outperformance Keeps Spreads Tight}
\end{figure}
\begin{figure}[htbp]
\centering
\includegraphics[width=0.88\linewidth]{figures/fig_002.jpg}
\caption{Exhibit 2 - Exhibit 2: Primary Balance and Real Growth are Close or Above the Historical Average}
\end{figure}
\begin{figure}[htbp]
\centering
\includegraphics[width=0.88\linewidth]{figures/fig_003.jpg}
\caption{Exhibit 3 - Exhibit 3: Lower Energy Exposure, Limited Fiscal Risk}
\end{figure}
\begin{figure}[htbp]
\centering
\includegraphics[width=0.88\linewidth]{figures/fig_004.jpg}
\caption{Exhibit 4 - Exhibit 4: Larger Energy Support, Strong Fiscal Position EMU3 includes Germany, France and Italy.}
\end{figure}
\begin{figure}[htbp]
\centering
\includegraphics[width=0.88\linewidth]{figures/fig_005.jpg}
\caption{Exhibit 4 - Exhibit 4: Larger Energy Support, Strong Fiscal Position EMU3 includes Germany, France and Italy.}
\end{figure}
\begin{figure}[htbp]
\centering
\includegraphics[width=0.88\linewidth]{figures/fig_006.jpg}
\caption{Exhibit 5 - Exhibit 5: Labour Reallocation Supports Productivity}
\end{figure}
\begin{figure}[htbp]
\centering
\includegraphics[width=0.88\linewidth]{figures/fig_007.jpg}
\caption{Exhibit 6 - Exhibit 6: Investment Recovery Gains Traction}
\end{figure}
\begin{figure}[htbp]
\centering
\includegraphics[width=0.88\linewidth]{figures/fig_008.jpg}
\caption{Exhibit 7 - Exhibit 7: Net Migration and Demographics Lift Housing Prices}
\end{figure}
\begin{figure}[htbp]
\centering
\includegraphics[width=0.88\linewidth]{figures/fig_009.jpg}
\caption{Exhibit 8 - Exhibit 8: Recovery Fund Support Fades Gradually European Recovery Fund: Total Disbursements and Spending}
\end{figure}
\begin{figure}[htbp]
\centering
\includegraphics[width=0.88\linewidth]{figures/fig_010.jpg}
\caption{Exhibit 9 - Exhibit 9: Spreads Price European, Not Domestic, Risk}
\end{figure}
\begin{figure}[htbp]
\centering
\includegraphics[width=0.88\linewidth]{figures/fig_014.jpg}
\caption{Exhibit 1 - Exhibit 1: Divergence between tech vs. non-tech production persists}
\end{figure}
\begin{figure}[htbp]
\centering
\includegraphics[width=0.88\linewidth]{figures/fig_015.jpg}
\caption{Exhibit 2 - Exhibit 2: Credit card sales pulled back in April}
\end{figure}
\begin{figure}[htbp]
\centering
\includegraphics[width=0.88\linewidth]{figures/fig_016.jpg}
\caption{Exhibit 3 - Exhibit 3: High-frequency data suggest inflation moderated in May, as energy price increases slowed and other goods prices fell Index (end-June, 2023=100)}
\end{figure}
\begin{figure}[htbp]
\centering
\includegraphics[width=0.88\linewidth]{figures/fig_017.jpg}
\caption{Exhibit 4 - Exhibit 4: Regular workers' wage growth continued to moderate, in contrast to the post-pandemic inflation cycle}
\end{figure}
\begin{figure}[htbp]
\centering
\includegraphics[width=0.88\linewidth]{figures/fig_018.jpg}
\caption{Exhibit 5 - Exhibit 5: Seoul housing prices have picked up again}
\end{figure}
\begin{figure}[htbp]
\centering
\includegraphics[width=0.88\linewidth]{figures/fig_019.jpg}
\caption{Exhibit 6 - Exhibit 6: Tech cycles have tended to move along with Seoul housing prices}
\end{figure}
{\small\color{refgray}\textbf{References:} [R001] 西班牙的结构性韧性被低估了，能源冲击只是压力测试; [R002] 韩国央行真正的转向，不是加息本身，而是对“结构性通胀”的承认}

\section{亚洲汇市：地缘风险缓释与人民币定价机制回归}
\textbf{美伊协议可能性上升将触发亚洲货币的系统性重估，人民币升值潜力被结构性低估；同时，人民币中间价模型显示逆周期因子收敛，定价机制正回归基准逻辑。}

\begin{itemize}
  \item 美伊协议若达成，市场低估的并非能源价格下行空间，而是地缘风险缓释后亚洲出口导向型货币的系统性估值修复，亚洲货币走势将回归各国基本面与资本流动博弈。
  \item 人民币升值潜力被多重因素共振低估：DXY下行时CNH敏感度高达54\%而上涨时仅24\%、中间价持续走低、企业结汇意愿强、中美关系改善，且人民币仍被低估约9.5\%。
  \item 人民币中间价模型预测值大幅下调503个基点至6.7870，逆周期因子正在系统性收敛，模型误差从-1800个基点峰值收窄至+400个基点，暗示央行干预力度减弱。
  \item 若逆周期因子收敛趋势延续，人民币汇率的双向波动将不再是口号，而是可交易的真实状态，对跨境资本流动和资产配置有不可忽视的含义。
\end{itemize}
{\small\color{refgray}\textbf{References:} [R003] 美伊协议或重塑亚洲汇市，最被低估的不是油价，而是人民币; [R004] 人民币中间价模型正在发出一个被忽视的定价信号}

\section{AI/算力与半导体：封装瓶颈重塑竞争格局}
\textbf{先进封装正从配套工艺演变为AI芯片性能与成本的决定性瓶颈，封装产能的稀缺性正在重新定义供给曲线和产业竞争格局。}

\begin{itemize}
  \item AI半导体TAM在2026年可能接近5000亿美元，全球云资本开支预计接近8000亿美元，但资本开支能否转化为实际算力取决于封装产能是否跟得上。
  \item CoWoS产能从2023年的15kwpm扩张到2027年的245kwpm，增长超16倍，但TSMC自身产能仅能覆盖约120kwpm，其余依赖外包，封装产能扩张速度直接决定AI芯片实际出货量。
  \item SoIC（系统级集成芯片）是下一个增长点，台积电产能规划从2024年的3kwpm拉至2028年的78kwpm，客户从英伟达和AMD扩展至苹果、高通、博通等。
  \item 当需求侧资本开支呈指数级上升而供给侧封装产能只能线性扩张时，封装从技术问题变成定价权和竞争格局的决定性变量，优先锁定产能的企业将获得出货保障。
\end{itemize}
{\small\color{refgray}\textbf{References:} [R008] AI半导体的下一个瓶颈不在芯片本身，而在封装}

\section{能源与大宗：安全冲击对供给侧的再定价}
\textbf{山西煤矿事故可能触发中央层面更严格的安全检查，市场低估了安全冲击对供给侧长期约束力的再定价能力，焦煤和动力煤短期均面临供给收缩压力。}

\begin{itemize}
  \item 山西煤矿事故导致沁源县全部25座在产煤矿停产自查，涉及年产能2740万吨炼焦煤，随后陕西、内蒙古、河南等地要求安全自查，波及产能可能超过8000万吨/年。
  \item 事故发生在传统用煤淡季，但淡季库存本就不高、市场对供给弹性缺乏预期，使得价格对供给收缩的反应比旺季更为敏感。
  \item 事故可能引发中央层面新一轮更严格的安全检查，进而影响动力煤供给，叠加迎峰度夏旺季，电厂可能加速补库，推动动力煤价格短期上行。
  \item 中国燃气行业面临成本端不对称风险：工业气量回暖具有假性复苏特征，驱动因素为钢铁行业强势和东南亚订单回流，而非内需触底反弹，可持续性存疑。
  \item 燃气行业单位毛差虽年初改善，但4月新合同年开启后成本压力已抬头，LNG现货价格高企压低接收站利用率，过去市场认为的第二增长曲线面临挑战。
\end{itemize}
{\small\color{refgray}\textbf{References:} [R006] 中国燃气行业：市场真正低估的不是需求，而是成本端的不对称风险; [R007] 市场低估的不是需求，而是安全冲击对供给侧的再定价能力}

\section{地产与消费：奢侈品品牌势能复苏与燃气接驳疲弱}
\textbf{奢侈品行业品牌势能复苏，美国市场为核心引擎，亚太企稳为信心指标；而中国燃气行业新接户数持续下滑，地产链相关需求依然低迷。}

\begin{itemize}
  \item 香奈儿FY25收入重回增长，美国市场以7\%有机增速领跑，销售额突破40亿美元创历史新高，亚太市场全年虽微降1\%但下半年已转正，全球需求底部已过。
  \item 香奈儿品牌势能复苏并非依赖单一爆款或设计师光环，而是品牌底层资产的重新激活，资本开支转向垂直整合，重点投向门店扩张和供应链。
  \item 中国燃气行业新接户数持续下滑：昆仑能源前四月新增约20万户（全年指引60-70万），华润燃气新接户数同比下降20\%，地产新开工疲软拖累接驳收入。
  \item 燃气行业工业气量回暖高度集中在钢铁等少数行业，与地产相关的玻璃和建材用气仍持续疲弱，并非全局性复苏。
\end{itemize}
{\small\color{refgray}\textbf{References:} [R006] 中国燃气行业：市场真正低估的不是需求，而是成本端的不对称风险; [R009] 香奈儿FY25财报的真正信号：品牌势能复苏比利润修复更值得关注}

\section{区域市场：港股定价权迁移与被动资金重构}
\textbf{恒生指数调整揭示被动资金正系统性地将港股定价重心从传统金融、能源和硬件制造转向互联网平台、医疗健康和AI软件服务。}

\begin{itemize}
  \item 此次恒指调整预计产生约110亿美元被动双向交易额，互联网/媒体娱乐板块净流入最大（+14亿美元），医疗健康紧随其后（+8.5亿美元）。
  \item 科技硬件/半导体、能源、银行、保险板块面临流出压力，科技硬件与半导体预计净流出约10亿美元，能源、银行和保险分别面临3.5亿至5亿美元流出。
  \item 恒生科技指数纳入MiniMax和智谱两家AI软件公司，同时剔除金山软件和金蝶国际，意味着科技指数正从硬件+软件混合配置向AI原生+平台型转变。
  \item 腾讯因重新上限调整获得近10亿美元被动买入，进一步拉高互联网板块权重，港股指数编制逻辑正从代表香港市场整体转向代表中国新经济未来。
\end{itemize}
\begin{figure}[htbp]
\centering
\includegraphics[width=0.88\linewidth]{figures/fig_026.jpg}
\caption{Exhibit 4 - Exhibit 5: HSI: Current additions initially underperformed but rebounded pre-announcement; historically, outperformance tends to continue before a reversal around the effective date}
\end{figure}
\begin{figure}[htbp]
\centering
\includegraphics[width=0.88\linewidth]{figures/fig_027.jpg}
\caption{Exhibit 6 - Exhibit 6: HSCEI: Current additions have sharply underperformed pre-announcement; historically, returns have moderated post-announcement before recovering after the effective date}
\end{figure}
\begin{figure}[htbp]
\centering
\includegraphics[width=0.88\linewidth]{figures/fig_028.jpg}
\caption{Exhibit 7 - Exhibit 7: HSTECH: Current additions have sharply outperformed; the outperformance typically continues from the announcement date through the effective date, then stabilizes}
\end{figure}
\begin{figure}[htbp]
\centering
\includegraphics[width=0.88\linewidth]{figures/fig_029.jpg}
\caption{Exhibit 8 - Exhibit 8: HSCI: Current additions are largely consistent with historical pre-announcement outperformance trends; historical patterns indicate further outperformance thereafter}
\end{figure}
\begin{figure}[htbp]
\centering
\includegraphics[width=0.88\linewidth]{figures/fig_030.jpg}
\caption{Exhibit 9 - Exhibit 9: Southbound Inclusions: SB ownership historical rose by 1pp within the first two days of inclusion, followed by 5pp over three months, while share prices often rally beforehand and decline post-inclusion before stabilizin}
\end{figure}
\begin{figure}[htbp]
\centering
\includegraphics[width=0.88\linewidth]{figures/fig_031.jpg}
\caption{Exhibit 10 - Exhibit 10: Southbound Exclusions: SB ownership typically drops by 1pp within the first two days of exclusion, followed by a further 3pp reduction over the next three months; share prices tend to face persistent pressure}
\end{figure}
{\small\color{refgray}\textbf{References:} [R005] 恒指调整的真正信号：被动资金正在重构港股定价权}

\section*{结语}
综上，市场主线正从单一通胀交易转向结构性分化：欧洲韧性被重估、亚洲汇市面临地缘风险缓释后的系统性重估、AI硬件瓶颈从芯片转向封装、安全冲击对供给侧再定价、以及港股定价权向新经济迁移。这些结构性变化为资产配置提供了新的分析框架。
\vfill
{\small\color{refgray}Personal reading notes and learning share only. Not investment advice.}
\end{document}
